\begin{textsample}{POS Dim 4 – 2015\_c5 – Score 1.00 – 2015\_c5/enade\_2015\_secr\_exec\_875f6cdd83.txt}  \label{ex:f4_pos_169}
CAPÍTULO 5 DISTRIBUIÇÃO DOS CONCEITOS 5.1 PANORAMA NACIONAL DA DISTRIBUIÇÃO DOS CONCEITOS O Conceito Enade é um indicador de qualidade que avalia o desempenho dos estudantes a partir dos resultados do Enade.
Os valores possíveis vão de 1 ( pior situação ) a 5 ( melhor situação ).
As Unidades de \textbf{Observação} com apenas um ou sem nenhum concluinte participante não obtêm o Conceito Enade, ficando Sem Conceito ( SC ).
A Tabela 5.1 apresenta a quantidade e distribuição de cursos de Secretariado Executivo participantes do Enade/2015, por faixa de conceito e Grande Região.
Mais uma vez, ressaltando que a diferença entre os cursos tabulados neste capítulo e no capítulo 2 correspondem aos cursos sem conceito, em princípio, aqueles sem alunos concluintes que participassem da prova ou que tivessem tão somente um aluno concluinte18.
Observando -se os dados da Tabela 5.1, nota -se que, dos 61 cursos participantes, 27 ( 44,3 \% ) classificaram -se com conceito 3, o valor modal.
Este foi também o conceito modal em mais quatro regiões : Nordeste ( 30,0 \% ), Sudeste ( 45,5 \% ), Sul ( 44,4 \% ) e Centro-Oeste ( 71,4 \% ).
Na região Norte, a moda foi o conceito 2 ( 75,0 \% ).
O conceito 2 foi o segundo mais frequente em nível nacional ( 18,0 \%, correspondendo a 11 cursos ) e o conceito 4, o terceiro ( 16,4 \%, correspondendo a dez cursos ).
Houve, ainda, três cursos ( 4,9 \% ) que receberam conceito 5 e três ( 4,9 \% ) cursos que receberam conceito 1.
Dos 61 cursos de Secretariado Executivo, sete cursos ( 11,5 \% ) ficaram sem conceito ( SC ). 18 Estes cursos com somente um concluinte não foram considerados no capítulo 2.
A região Norte participou com quatro cursos ou 6,6 \% do total nacional.
Nesta região, o conceito 2 foi atribuído a três cursos dos quatro participantes, o que equivale a 75,0 \% do total regional, além de ser o conceito modal na região, e o conceito 3 a um curso restante.
Nenhum curso ficou sem conceito ( SC ).
A região Nordeste participou com dez cursos ou 16,4 \% do total nacional.
Nesta região, três cursos, 30,0 \% em termos regionais, obtiveram conceito 3, o conceito modal para a região.
Os conceitos 2, 4 e 5 foram atribuídos a um curso cada ( 10,0 \% cada ).
Dois cursos ( 20,0 \% ) obtiveram o conceito 1, e outros dois cursos ficaram sem conceito ( SC ).
Dos 22 cursos participantes da região Sudeste ( 36,1 \% do total nacional ), o conceito 3 também foi o conceito modal, tendo sido obtido por dez cursos ( 45,5 \% ).
Os conceitos 1 e 5 foram atribuídos a um curso cada ( 4,5 \% cada ), o conceito 2 a três cursos ( 13,6 \% ) e o conceito 4 a cinco cursos ( 22,7 \% ).
Os dois cursos restantes ( 9,1 \% ) ficaram sem conceito ( SC ).
Dos 18 cursos da região Sul, observa -se que há predominância do conceito 3 com 44,4 \% ( conceito modal ), ou seja, oito dos 18 cursos participantes na região Sul.
O conceito 2 foi atribuído a dois cursos ( 11,1 \% ) e o conceito 4, a quatro cursos ( 22,2 \% ).
O conceito 5 foi recebido por apenas um curso ( 5,6 \% ).
Nenhum curso na região Sul foi avaliado com o conceito 1 e três cursos ( 16,7 \% ) ficaram sem conceito ( SC ).
Já dos sete cursos participantes na região Centro-Oeste, cinco receberam conceito 3, o conceito modal ( 71,4 \% ) e dois ( 28,6 \% ) o conceito 2.
Nesta região, nenhum curso ficou sem conceito ( SC ) ou obteve os demais conceitos.
O Gráfico 5.1 apresenta a distribuição do Conceito Enade segundo Grande Região : poligonais mais à esquerda representam Grandes Regiões com pior distribuição, e poligonais mais à direita, Grandes Regiões com melhores conceitos.
A região Norte ( linha azul ) apresenta os cursos em IES com conceitos mais baixos, ao passo que as regiões Sudeste ( linha vermelha ) e Sul ( linha roxa ) apresentam as distribuições com valores maiores.
Cursos em IES das regiões Nordeste ( linha verde ) e Centro-Oeste ( linha preta ) apresentam situação intermediária. 5.2 CONCEITOS POR CATEGORIA ADMINISTRATIVA E POR GRANDE REGIÃO A Tabela 5.2 apresenta a distribuição dos cursos participantes do Enade/2015 de Secretariado Executivo, por Categoria Administrativa, de acordo com os conceitos por eles alcançados, segundo as Grandes Regiões brasileiras.
Dos 61 cursos participantes, 18 ( 29,5 \% ) eram ministrados em Instituições Públicas e 43 ( 70,5 \% ), em Instituições Privadas.
De acordo com as informações da Tabela 5.2, em termos nacionais, observa -se que todas as três IES que receberam o conceito 5, eram Públicas.
Dos 18 cursos participantes de IES Públicas, o conceito 3 foi o valor modal, atribuído a cinco cursos ( 27,8 \% ).
Entre os demais cursos de Instituições Públicas participantes, um curso ( 5,6 \% ) foi avaliado com conceito 1, oito cursos obtiveram os conceitos 2 e 4, quatro cursos cada ( 22,2 \% cada ) e, como já comentado, três foram alocados no conceito 5 ( 16,7 \% ).
Nesta categoria, apenas um curso ficou sem conceito ( SC ).
Na rede privada, o conceito modal foi 3, atribuído a 22 cursos dos 43 da categoria.
Entre os demais cursos participantes, dois receberam conceito 1, e sete, conceito 2.
O conceito 4 foi atribuído a seis cursos, e nenhum curso foi avaliado com o conceito 5.
Nesta Categoria Administrativa, seis cursos ficaram sem conceito ( SC ).
Na análise por região, observa -se que, na região Norte, as Instituições Privadas participaram com dois cursos ( 50,0 \% do total regional ), dos quais nenhum ficou sem conceito ( SC ) ou com conceitos 1, 4 ou 5.
Um curso foi avaliado com conceito 2 e o outro, com conceito 3, correspondendo a 50,0 \% cada.
As Instituições Públicas participaram com dois cursos na região Norte ( 50,0 \% do total regional ), os quais obtiveram conceito 2.
Desse modo, nenhum curso ficou sem conceito ( SC ).
Na região Nordeste, a rede privada concentrou quatro dos dez cursos participantes, o equivalente a 40,0 \% do total da região, a menor das proporções desta categoria entre todas as regiões.
Dentre os cursos oferecidos por IES Privadas no Nordeste, um curso ficou alocado no conceito 1 e um outro curso obteve conceito 3.
Dois cursos ficaram sem conceito ( SC ) enquanto nenhum recebeu os demais conceitos.
As Instituições Públicas dessa região participaram com seis cursos ( 60,0 \% ), dos quais dois ( 33,3 \% ) obtiveram conceito 3, o valor modal, e os quatro cursos restante receberam os demais conceitos, um curso cada ( 16,7 \% cada ).
Nenhum curso ficou sem conceito ( SC ) nesta combinação de categoria e região.
Na região Sudeste, a proporção de cursos da rede privada, 95,5 \%, foi mais elevada do que nas demais regiões brasileiras, correspondendo a 21 dos 22 cursos participantes.
Nesta categoria, na região Sudeste, o conceito modal foi 3 ( dez cursos ).
Nesta combinação de Categoria Administrativa e Grande Região, dois cursos não receberam conceito ( SC ) e nenhum foi avaliado com o conceito 5.
Os demais foram avaliados com conceito 1 ( um curso ), conceito 2 ( três cursos ) e conceito 4 ( cinco cursos ).
Há apenas um curso ( 4,5 \% ) oferecido por Instituição Pública na região Sudeste, sendo que este recebeu o conceito cinco.
As Instituições Privadas concentraram dez dos 18 cursos participantes da região Sul, 55,6 \% do total regional.
Desses, seis ficaram com conceito 3, o conceito modal.
Nesta combinação de Categoria Administrativa e Grande Região, dois cursos ficaram com conceitos 2 e 4, um curso cada.
Dois cursos ficaram sem conceito ( SC ) e nenhum recebeu conceito 1 ou 5.
As Instituições Públicas na região Sul participaram com oito cursos ( 44,4 \% ), e o conceito modal foi 4, atribuído a três cursos.
Dois cursos foram avaliados com conceitos 2 e 5 ( um curso, cada ) e dois cursos obtiveram conceito 3.
Apenas um curso ficou sem conceito ( SC ) enquanto nenhum recebeu o conceito 1.
Na região Centro-Oeste, seis dos sete cursos participantes eram de Instituições Privadas ( 85,7 \% em termos regionais ).
Destes, quatro concentraram -se no conceito 3, conceito modal e os outros dois receberam conceito 2.
Nenhum curso ficou sem conceito ( SC ) ou recebeu os demais conceitos nesta categoria.
Há apenas um curso

% matched lemmas: observação
\end{textsample}
